\subsection{Linear Classification and Generalized Linear Models}
\begin{definition} \label{def:sl-discriminant-functions-linear-boundaries} \hlt{Discriminant Functions and Linear Boundaries}\\
For $\mathcal{G}=\{1,\ldots,K\}$ and scores $\delta_k:\mathcal{X}\to\R$, set $\widehat{G}(x)\in\operatorname*{argmax}_{k\in\mathcal{G}}\delta_k(x)$.  Only score differences matter.
\begin{enumerate}[label=(\roman*)]
\setlength{\itemsep}{0pt}
\item \hlt{Pairwise boundary}: The $k$ versus $\ell$ boundary lies in $\{x:\delta_k(x)=\delta_\ell(x)\}$.
\item \hlt{Linear boundary}: If $\delta_k-\delta_\ell$ is affine in $x$ for all $k,\ell$, then all pairwise decision boundaries are hyperplanes.
\item \hlt{Binary log-odds}: If $\log\frac{\mathbb{P}(G=1\mid X=x)}{\mathbb{P}(G=2\mid X=x)}=\beta_0+\beta^\top x$, then the boundary is $\{x:\beta_0+\beta^\top x=0\}$.
\end{enumerate}
\end{definition}

\begin{method} \label{meth:sl-indicator-matrix-regression-classifier} \hlt{Indicator-Matrix Regression Classifier}\\
For classes $g_i\in\{1,\ldots,K\}$, set $t_k=e_k\in\R^K$, $Y_i=t_{g_i}$, and $\widetilde{x}_i^\top=(1,x_i^\top)$, coding labels as Euclidean targets.
\begin{enumerate}[label=(\roman*)]
\setlength{\itemsep}{0pt}
\item \hlt{Least-squares fit}: Choose $\widehat{B}\in\operatorname*{argmin}_B\sum\limits_{i=1}^{n}\norm{Y_i-B^\top\widetilde{x}_i}_2^2$; if $\widetilde{X}^\top\widetilde{X}$ is nonsingular, $\widehat{B}=(\widetilde{X}^\top\widetilde{X})^{-1}\widetilde{X}^\top Y$.
\item \hlt{Scores}: For a new $x$, set $\widetilde{x}=(1,x^\top)^\top$ and $\widehat{f}(x)\defeq\widehat{B}^\top\widetilde{x}$.
\item \hlt{Classification}: Predict $\widehat{G}(x)\in\operatorname*{argmax}_k\widehat{f}_k(x)$, equivalently nearest $t_k$.
\end{enumerate}
With an intercept, $\sum_k\widehat{f}_k(x)=1$, but $\widehat{f}_k(x)$ may leave $[0,1]$ and, for $K\geq3$, \hlt{mask} middle classes.
\end{method}

\begin{definition} \label{def:sl-generalized-linear-model} \hlt{Generalized Linear Model}\\
A \hlt{generalized linear model} gives $Y\mid X=x$ an exponential-family law and makes a linked mean linear:\\
$g(\mu(x))=w^\top x+b$. Logistic regression uses the Bernoulli law and logit link $g(\mu)=\log\frac{\mu}{1-\mu}$.
\end{definition}

Thus a GLM separates score, mean scale, and likelihood.

\begin{definition} \label{def:sl-multiclass-logistic-regression} \hlt{Multiclass Logistic Regression}\\
With class $K$ as baseline, set $\eta_k(x)\defeq\beta_{k0}+\beta_k^\top x=\log\frac{\mathbb{P}(G=k\mid X=x)}{\mathbb{P}(G=K\mid X=x)}$ for $k<K$.
\begin{enumerate}[label=(\roman*)]
\setlength{\itemsep}{0pt}
\item \hlt{Softmax probabilities}: For $k<K$, $p_k(x;\theta)\defeq\frac{\exp(\eta_k(x))}{1+\sum\limits_{\ell=1}^{K-1}\exp(\eta_\ell(x))}$, and $p_K(x;\theta)\defeq\frac{1}{1+\sum\limits_{\ell=1}^{K-1}\exp(\eta_\ell(x))}$.
\item \hlt{Likelihood}: Fit $\theta$ by maximising $\ell(\theta)\defeq\sum\limits_{i=1}^{n}\log p_{g_i}(x_i;\theta)$.
\end{enumerate}
The softmax gives $p(x;\theta)\in\Delta^{K-1}$; the reference class changes only the parametrisation.
\end{definition}

\begin{method} \label{meth:sl-logistic-irls} \hlt{Binary Logistic Regression by IRLS}\\
For $y_i\in\{0,1\}$ and intercept-augmented rows $x_i^\top$, set $p_i(\beta)\defeq(1+\exp(-x_i^\top\beta))^{-1}$.
\begin{enumerate}[label=(\roman*)]
\setlength{\itemsep}{0pt}
\item \hlt{Derivatives}: With $p=(p_i)$ and $W\defeq\operatorname{diag}(p_i(1-p_i))$, $\nabla\ell(\beta)=X^\top(y-p)$ and $\nabla^2\ell(\beta)=-X^\top W X$.
\item \hlt{Adjusted response}: At $\beta^{\mathrm{old}}$, set $z\defeq X\beta^{\mathrm{old}}+W^{-1}(y-p)$.
\item \hlt{IRLS update}: Set $\beta^{\mathrm{new}}\defeq(X^\top W X)^{-1}X^\top Wz$.
\end{enumerate}
IRLS is Newton's method as weighted least squares.\\
The log-likelihood is concave when the MLE is finite; step-halving controls Newton overshoot.
\end{method}

\begin{method} \label{meth:sl-l1-logistic-regression} \hlt{$L_1$-Regularized Logistic Regression}\\
With unpenalized intercept, maximize $\ell_\lambda(\beta_0,\beta)\defeq\sum\limits_{i=1}^{n}\{y_i\eta_i-\log(1+\exp(\eta_i))\}-\lambda\sum\limits_{j=1}^{p}\abs{\beta_j}$, where $\eta_i\defeq\beta_0+x_i^\top\beta$.
\begin{enumerate}[label=(\roman*)]
\setlength{\itemsep}{0pt}
\item \hlt{Probabilities}: Set $p_i\defeq(1+\exp(-\eta_i))^{-1}$ and residual $r\defeq y-p$.
\item \hlt{KKT conditions}: If $\widehat{\beta}_j\neq0$, then $x_j^\top r=\lambda\sign(\widehat{\beta}_j)$; if $\widehat{\beta}_j=0$, then $\abs{x_j^\top r}\leq\lambda$.
\end{enumerate}
The KKT conditions are thresholded feature--residual correlations.
\end{method}
